% БҮЛЭГ 3 — СИСТЕМИЙН АРХИТЕКТУР БА ЗАГВАРЫН ЗОХИОН БАЙГУУЛАЛТ

\chapter{Системийн архитектур ба загварын зохион байгуулалт}
\label{Chapter3}

Энэхүү бүлэгт судалгааны ажлын хүрээнд дэвшүүлсэн туршилтын архитектур, загварын
сургалтын pipeline, интерфейсийн бүтэц, болон үнэлгээний арга зүйг нарийвчлан
тайлбарлана. Бидний үндсэн таамаглал (hypothesis) нь хоёр шатлалт ангиллын
архитектур ашиглан ``Negative'' агуулгыг нарийвчлан ялгах замаар системийн
ерөнхий нарийвчлалыг нэмэгдүүлэх байсан бөгөөд уг бүтцийг дараах байдлаар
зохион байгуулсан.

%-------------------------------------------------------------------------------
\section{Системийн туршилтын архитектур}

Системийн зохион байгуулалтыг илүү тодорхой болгохын тулд эхлээд системийн
хэрэглээний тохиолдлын (Use Case) болон үйл ажиллагааны (Activity) логикийг
тодорхойлов.

\subsection{Хэрэглээний тохиолдлын диаграмм (Use Case Diagram)}

Зураг \ref{fig:use-case}-т харуулсанчлан систем нь гурван үндсэн оролцогчтой:
Өсвөр насны хэрэглэгч, Модератор (зохицуулагч) болон Системийн администратор.
Хэрэглэгч сэтгэгдэл оруулах бол, Модератор нь системийн ангилсан үр дүнг хянаж,
хортой агуулгыг шүүх, бүтээлч шүүмжлэлийг зөвшөөрөх үйлдлийг гүйцэтгэнэ.

\begin{figure}[H]
    \centering
    \begin{tikzpicture}[
        node distance=1.2cm and 2cm,
        usecase/.style={ellipse, draw, minimum width=3.5cm, minimum height=1.1cm, text centered, font=\small, fill=blue!5, drop shadow, align=center},
        actor/.style={circle, draw, minimum size=1cm, fill=orange!10, drop shadow},
        line/.style={draw, -latex, thick}
    ]
        % System Boundary
        \draw[thick, dashed, fill=gray!2] (-3, -4.5) rectangle (3, 4.5);
        \node at (0, 4.1) {\textbf{NLP Систем}};

        % Use Cases
        \node (uc1) [usecase] at (0, 3) {Сэтгэгдэл оруулах};
        \node (uc2) [usecase] [below=of uc1] {Ангиллын дүн үзэх};
        \node (uc3) [usecase] [below=of uc2] {Хортойг шүүх};
        \node (uc4) [usecase] [below=of uc3] {Бүтээлчийг зөвшөөрөх};
        \node (uc5) [usecase] [below=of uc4] {Загварыг сургах};

        % Actors
        \node (admin) [actor, left=4cm of uc3] {};
        \node [below=0.2cm of admin] {\small Администратор};
        
        \node (user) [actor, right=4cm of uc1] {};
        \node [below=0.2cm of user] {\small Хэрэглэгч};

        % Relations
        \draw [line] (user) -- (uc1);
        \draw [line] (admin) -- (uc2.west);
        \draw [line] (admin) -- (uc3.west);
        \draw [line] (admin) -- (uc4.west);
        \draw [line] (admin) -- (uc5.west);
    \end{tikzpicture}
    \caption{Системийн хэрэглээний тохиолдлын диаграмм}
    \label{fig:use-case}
\end{figure}

\subsection{Үйл ажиллагааны диаграмм (Activity Diagram)}

Системийн логик урсгалыг Зураг \ref{fig:activity-diagram}-т харуулав. Энэхүү
диаграмм нь хоёр шатлалт ангиллын шийдвэр гаргах үйл явцыг нарийвчлан харуулж
байна.

\begin{figure}[H]
    \centering
    \begin{tikzpicture}[
        node distance=1cm and 1.5cm,
        start/.style={circle, draw, fill=black, inner sep=3pt},
        stop/.style={circle, draw, double, fill=black!50, inner sep=2.5pt},
        procstep/.style={rectangle, draw, rounded corners, minimum width=4cm, minimum height=1cm, text centered, fill=green!5, drop shadow, align=center},
        decision/.style={diamond, draw, aspect=2, minimum width=3cm, text centered, fill=yellow!10, font=\small, drop shadow, inner sep=0pt, align=center},
        arrow/.style={thick, ->, >=stealth}
    ]
        \node (init) [start] {};
        \node (input) [procstep, below=of init] {Текст хүлээн авах};
        \node (prep) [procstep, below=of input] {Урьдчилсан боловсруулалт};
        \node (s1) [decision, below=1.5cm of prep] {Сөрөг? (Stage 1)};
        
        \node (allow) [procstep, left=2.5cm of s1] {Зөвшөөрөх (Allow)};
        
        \node (s2) [decision, below=1.5cm of s1] {Хортой? (Stage 2)};
        \node (block) [procstep, right=2.5cm of s2] {Блок хийх (Block)};
        
        \node (end) [stop, below=of s2] {};

        \draw [arrow] (init) -- (input);
        \draw [arrow] (input) -- (prep);
        \draw [arrow] (prep) -- (s1);
        \draw [arrow] (s1) -- node[anchor=east, xshift=-0.1cm] {Тийм} (s2);
        \draw [arrow] (s1) -- node[anchor=south] {Үгүй} (allow);
        \draw [arrow] (s2.west) -| node[anchor=south, xshift=1.5cm] {Үгүй} (allow);
        \draw [arrow] (s2) -- node[anchor=south] {Тийм} (block);
        
        \draw [arrow] (allow) |- (end);
        \draw [arrow] (block) |- (end);
    \end{tikzpicture}
    \caption{Системийн шүүлтүүрийн логикийн үйл ажиллагааны диаграмм}
    \label{fig:activity-diagram}
\end{figure}

\subsection{Давхаргын архитектур}

Системийн архитектур нь гурван үндсэн давхаргаас бүрдэнэ: (1)~өгөгдлийн
боловсруулалтын давхарга (Data Processing Layer), (2)~загварын давхарга (Model
Layer), (3)~хэрэглэгчийн интерфейсийн давхарга (Presentation Layer). Эдгээр
давхарга нь бие биенээс харьцангуй бие даасан (loosely coupled) бүтэцтэй бөгөөд
давхарга бүрийг тус тусад нь шинэчлэх, туршихад хялбар.

\begin{figure}[H]
  \centering
  \begin{tikzpicture}[
      block/.style={rectangle, draw, fill=blue!10, text width=10em,
                    text centered, rounded corners, minimum height=3em, align=center},
      line/.style={draw, -latex}
    ]
    \node [block] (input) {Хам текст \\ (news.mn, Facebook)};
    \node [block, below of=input, node distance=3cm] (preprocess)
          {Урьдчилсан боловсруулалт \\ (Цэвэрлэгээ, Токенчлол)};
    \node [block, below of=preprocess, node distance=3cm] (stage1)
          {Stage~1: Сентимент ангилал \\ (Positive / Neutral / Negative)};
    \node [block, below of=stage1, node distance=3cm] (stage2)
          {Stage~2: Нарийвчилсан ангилал \\ (Toxic vs Constructive)};
    \node [block, below of=stage2, node distance=3cm] (output)
          {Үр дүн \\ (Gradio интерфейс)};
    \path [line] (input) -- (preprocess);
    \path [line] (preprocess) -- (stage1);
    \path [line] (stage1) -- node [right, xshift=0.3cm]
          {\small Negative} (stage2);
    \path [line] (stage2) -- (output);
  \end{tikzpicture}
  \caption{Системийн ерөнхий архитектурын схем}
  \label{fig:system-arch}
\end{figure}

Зураг~\ref{fig:system-arch}-т үзүүлсэн ерөнхий бүтцийн дагуу хэрэглэгчийн
оруулсан текст эхлээд урьдчилсан боловсруулалтын шатыг давж, дараа нь хоёр
дараалсан ангиллын шатаар дамжина. Stage~1 нь текстийг Positive, Neutral,
Negative гэсэн гурван ангилалд хуваана. Зөвхөн Negative гэж ангилагдсан
текст Stage~2 руу дамжих бөгөөд тэнд Toxic (хортой) ба Constructive
(бүтээлч шүүмжлэл) гэсэн хоёр ангилалд ялгана.

%-------------------------------------------------------------------------------
\section{Хоёр шатлалт ангиллын архитектур}

Хоёр шатлалт (two-stage) ангиллын архитектурыг сонгосон гол шалтгаан нь дангаар
нэг загвараар дөрвөн ангилалт (Positive, Neutral, Toxic, Constructive) хийхэд
ангиллуудын хоорондын семантик ялгаа тэнцвэргүй байдаг явдал юм. Positive болон
Neutral нь Negative-ээс харьцангуй амархан ялгагддаг бол Toxic болон
Constructive хоёрын хоорондын хил маш нарийн. Тиймд pipeline-ийг хоёр
тусдаа шатанд хуваах нь загвар бүрийг өөрийн тусгай даалгаварт
(specialized task) төвлөрүүлэх боломж олгоно.

\subsection{Stage~1 --- Сентимент ангилал}

Stage~1 загварын зорилго нь оролтын текстийг Positive, Neutral, Negative гэсэн
гурван ангилалд хуваах юм. Энэ шатанд дараах загваруудыг хэрэгжүүлж,
харьцуулна:

\begin{itemize}
  \item \textbf{MN-BERT (fine-tuned):} Монгол хэлний урьдчилсан сургалттай
        (pre-trained) BERT загварыг гурван ангиллын dataset дээр fine-tuning
        хийнэ. Сүүлийн \texttt{[CLS]} токены давхаргын гаралтыг softmax
        ангилагч руу дамжуулна.
  \item \textbf{Logistic Regression (LR):} TF-IDF шинж чанарын дээр L2
        регуляризацитай LR загвар сургана.
  \item \textbf{Naive Bayes (NB):} TF-IDF шинж чанарын дээр Multinomial
        Naive Bayes ашиглана.
  \item \textbf{SVM:} Linear kernel бүхий Support Vector Machine загвар
        сургана.
\end{itemize}

Stage~1-ийн гол шаардлага нь Negative ангилалын recall-ийг аль болох
өндөр байлгах явдал юм --- учир нь Negative текстийг алдах (False Negative)
тохиолдолд тухайн текст Stage~2 руу шилжихгүй, Toxic агуулга илрүүлэгдэхгүй.

\subsection{Stage~2 --- Toxic vs Constructive ялгалт}

Stage~1-ээс Negative гэж ангилагдсан текстүүд Stage~2 руу дамжина. Энэ шатанд
хоёр ангилалт (binary classification) хийнэ:

\begin{itemize}
  \item \textbf{Toxic:} Доромжлол, хувь хүнд чиглэсэн довтолгоо, үзэн ядалтын
        үг, кибер дарамт зэрэг устгах шаардлагатай агуулга.
  \item \textbf{Constructive Criticism:} Аргумент, нотолгоонд суурилсан сөрөг
        үнэлгээ, бодлогыг шүүмжлэх, асуудлыг тодорхойлох зэрэг хамгаалах
        шаардлагатай агуулга.
\end{itemize}

Stage~2-т мөн MN-BERT fine-tuned загвар болон суурь загваруудыг (LR, NB, SVM)
хэрэгжүүлж харьцуулна. Энэ шатны тусгай бэрхшээл нь Toxic болон Constructive
хоёрын хоорондын семантик ойролцоо байдал бөгөөд хоёулаа сөрөг хандлагыг
агуулдаг боловч илэрхийлэх хэлбэр, зорилго нь ялгаатай.

%-------------------------------------------------------------------------------
\section{MN-BERT загварын fine-tuning}

MN-BERT нь Монгол хэлний кирилл корпус дээр урьдчилан сургагдсан (pre-trained)
BERT загвар юм. Fine-tuning процесст дараах тохиргоонуудыг ашиглана.

\subsection{Загварын гиперпараметрүүд}

\begin{table}[H]
\centering
\caption{MN-BERT fine-tuning гиперпараметрүүд}
\label{tab:hyperparams}
\begin{tabular}{@{}ll@{}}
\toprule
\textbf{Параметр} & \textbf{Утга} \\
\midrule
Learning rate             & $2 \times 10^{-5}$ \\
Batch size                & 16 \\
Epoch тоо                 & 3--5 (early stopping) \\
Max sequence length       & 128 токен \\
Optimizer                 & AdamW \\
Weight decay              & 0.01 \\
Warmup steps              & Нийт алхмын 10\% \\
Loss function (Stage~1)   & Cross-entropy \\
Loss function (Stage~2)   & Binary cross-entropy \\
\bottomrule
\end{tabular}
\end{table}

\subsection{Fine-tuning стратеги}

Fine-tuning нь хоёр шатанд тус тусдаа хийгдэнэ. Stage~1-ийн загвар нь бүх
dataset (Positive, Neutral, Negative) дээр гурван ангиллын cross-entropy
loss-оор сургагдана. Stage~2-ийн загвар нь зөвхөн Negative шошготой
дэд-dataset дээр binary cross-entropy loss-оор сургагдана.

Хоёр загвар хоёулаа MN-BERT-ийн ижил урьдчилсан сургалтын жинг (pre-trained
weights) эхний цэг болгон ашигладаг боловч fine-tuning-ийн дараа тус тусдаа
тохируулагдсан (specialized) загварууд үүснэ.

\textbf{Early stopping.} Validation loss 2 epoch дараалсан буурахгүй бол
сургалтыг зогсооно. Энэ нь хэт тохируулга (overfitting)-аас сэргийлэх
стандарт арга юм.

\textbf{Class-weighted loss.} Ангиллуудын тэнцвэргүй тохиолдолд loss
function-д class weight нэмж, цөөнх ангилалд (minority class) илүү өндөр
торгууль (penalty) тооцоолно. Weight нь ангилал бүрийн жишээний
урвуу давтамжаас тооцогдоно.

%-------------------------------------------------------------------------------
\section{Суурь загварууд (Baseline Models)}

MN-BERT-ийн гүйцэтгэлийг үнэлэхийн тулд гурван суурь загвартай харьцуулна.
Суурь загварууд бүгд ижил preprocessing pipeline-аар боловсруулагдсан
өгөгдлийн дээр сургагдана.

\begin{table}[H]
\centering
\caption{Суурь загваруудын тохиргоо}
\label{tab:baselines}
\begin{tabular}{@{}lll@{}}
\toprule
\textbf{Загвар} & \textbf{Шинж чанар} & \textbf{Тохиргоо} \\
\midrule
Logistic Regression & TF-IDF (unigram + bigram) & L2 регуляризаци, $C=1.0$ \\
Naive Bayes         & TF-IDF (unigram + bigram) & Multinomial NB \\
SVM                 & TF-IDF (unigram + bigram) & Linear kernel, $C=1.0$ \\
\bottomrule
\end{tabular}
\end{table}

TF-IDF (Term Frequency--Inverse Document Frequency) нь текстийн статистик
шинж чанарыг тоон вектор хэлбэрт хөрвүүлэх арга бөгөөд unigram болон bigram
хослолыг ашиглана. Эдгээр суурь загварууд нь MN-BERT-ийн нэмүү үнэ цэнийг
(value-add) тодорхойлоход туслах benchmark болно.

%-------------------------------------------------------------------------------
\section{Тэнцвэргүй өгөгдлийн шийдэл}

Бодит орчны өгөгдөлд ангиллуудын тоон тэнцвэр хангагдах нь ховор. Тухайлбал,
Toxic агуулга нь Constructive-ээс тоогоор хамаагүй олон, Neutral нь Positive-ээс
давамгайлах зэрэг тэнцвэргүй байдал гарч болно. Энэ асуудлыг шийдвэрлэхэд
дараах аргуудыг хэрэгжүүлнэ:

\begin{itemize}
  \item \textbf{SMOTE (Synthetic Minority Over-sampling Technique):} Суурь
        загваруудын (LR, NB, SVM) TF-IDF шинж чанарын дээр SMOTE ашиглан
        цөөнх ангилалд хиймэл жишээ (synthetic sample) үүсгэнэ.
  \item \textbf{Class-weighted loss:} MN-BERT загварын loss function-д
        ангилал бүрийн урвуу давтамжид суурилсан жин (class weight) нэмнэ.
  \item \textbf{Stratified sampling:} Dataset хуваалт (train/val/test) дээр
        ангилал бүрийн харьцааг тэнцүү хадгална.
\end{itemize}

%-------------------------------------------------------------------------------
\section{Gradio хэрэглэгчийн интерфейс}

Системийн хэрэглэгчийн интерфейсийг Gradio framework ашиглан хэрэгжүүлнэ.
Gradio нь Python хэл дээр суурилсан, хурдан прототип бүтээхэд тохиромжтой
вэб framework юм. Интерфейсийн гол функцууд:

\begin{enumerate}
  \item \textbf{Текст оруулах:} Хэрэглэгч нэг текст эсвэл олон текст (batch)
        оруулж ангилуулах боломжтой.
  \item \textbf{Ангиллын үр дүн харуулах:} Stage~1 болон Stage~2-ийн
        ангиллын үр дүнг итгэлцлийн хувь (confidence score)-тай хамт харуулна.
  \item \textbf{Дүрслэл (visualization):} Ангиллын тархалтын график,
        confidence distribution зэргийг дүрсэлж харуулна.
  \item \textbf{Batch шинжилгээ:} CSV файл хэлбэрээр олон текст нэг дор
        оруулж, үр дүнг татаж авах (download) боломжтой.
\end{enumerate}

\begin{figure}[H]
  \centering
  \begin{tikzpicture}[
      box/.style={rectangle, draw, fill=green!10, text width=8em,
                  text centered, rounded corners, minimum height=2.5em, align=center},
      line/.style={draw, -latex}
    ]
    \node [box] (input) {Текст оруулах};
    \node [box, right of=input, node distance=5cm] (api) {Загвар дуудах \\ (MN-BERT API)};
    \node [box, below of=api, node distance=3cm] (result) {Үр дүн харуулах};
    \node [box, left of=result, node distance=5cm] (viz) {Дүрслэл \\ (График, хүснэгт)};
    \path [line] (input) -- (api);
    \path [line] (api) -- (result);
    \path [line] (result) -- (viz);
  \end{tikzpicture}
  \caption{Gradio интерфейсийн ажиллагааны бүтэц}
  \label{fig:streamlit-flow}
\end{figure}

%-------------------------------------------------------------------------------
\section{Үнэлгээний арга зүй}

Загваруудын гүйцэтгэлийг дараах стандарт NLP үнэлгээний үзүүлэлтүүдээр хэмжинэ.

\subsection{Ангиллын үзүүлэлтүүд}

\begin{itemize}
  \item \textbf{Accuracy} --- нийт зөв ангилагдсан жишээний хувь.
  \item \textbf{Precision} --- тухайн ангилалд хамааруулсан жишээнүүдээс хэдэн
        хувь нь үнэхээр тэр ангилалд хамаарах.
  \item \textbf{Recall} --- тухайн ангилалд бодитоор хамаарах жишээнүүдээс хэдэн
        хувь нь зөв илрүүлэгдсэн.
  \item \textbf{$F_1$-score} --- Precision болон Recall-ийн гармоник дундаж.
        Тэнцвэргүй dataset-ийн нөхцөлд Accuracy-ээс илүү бодитой үнэлгээ өгдөг.
  \item \textbf{Confusion Matrix} --- ангилал бүрийн зөв болон буруу
        таамаглалын матриц.
  \item \textbf{AUC-ROC} --- загварын ялгах чадварыг босгоноос үл хамааран
        хэмжих муруйн доорх талбай.
\end{itemize}

\subsection{Харьцуулалтын арга зүй}

Загвар бүрийг ижил test set дээр үнэлнэ. Харьцуулалтыг шударгаар хийхийн тулд:
\begin{enumerate}
  \item Бүх загвар ижил train/val/test хуваалтыг ашиглана.
  \item Бүх загвар ижил preprocessing pipeline-ийг ашиглана (MN-BERT-ийн хувьд
        tokenizer нь өөрийнх, суурь загваруудын хувьд TF-IDF).
  \item Test set нь сургалтын явцад огт ашиглагдаагүй --- зөвхөн эцсийн
        үнэлгээнд нэг удаа ашиглагдана.
  \item Macro-average $F_1$-score-ийг гол харьцуулалтын үзүүлэлт болгоно.
\end{enumerate}

%-------------------------------------------------------------------------------
\section{Технологийн стек}

Системийг хэрэгжүүлэхэд ашиглагдах гол технологиудыг хүснэгт~\ref{tab:tech-stack}-д нэгтгэв.

\begin{table}[H]
\centering
\caption{Системийн технологийн стек}
\label{tab:tech-stack}
\begin{tabular}{@{}ll@{}}
\toprule
\textbf{Зориулалт} & \textbf{Технологи} \\
\midrule
Програмчлалын хэл       & Python 3.10+ \\
Загварын framework       & PyTorch, Hugging Face Transformers \\
Урьдчилсан загвар       & MN-BERT (Hugging Face Hub) \\
Суурь загварууд          & scikit-learn \cite{pedregosa2011} \\
Тэнцвэргүй өгөгдөл      & imbalanced-learn (SMOTE) \\
Хэрэглэгчийн интерфейс  & Gradio \\
Өгөгдлийн боловсруулалт  & pandas, NumPy \\
Дүрслэл                  & matplotlib, seaborn \\
Хувилбар хяналт          & Git, GitHub \\
\bottomrule
\end{tabular}
\end{table}

%-------------------------------------------------------------------------------
\section{AI хэрэгслийн ашиглалт}
\label{sec:ai-disclosure}

Guideline-ийн зүйл 7-д заасны дагуу дипломын ажилд ашигласан AI хэрэгслийг
ил тодоор мэдүүлж байна. Ашигласан гол хэрэгслүүд: \textbf{Anthropic Claude}
(Sonnet/Opus хувилбарууд) ба \textbf{Google Gemini}.

\begin{itemize}
  \item \textbf{Кодын хөгжүүлэлтэд (Claude, Gemini):} PyTorch, HuggingFace
    Transformers, \texttt{imbalanced-learn}, Gradio зэрэг нээлттэй эх
    сурвалжтай сангуудын албан ёсны баримт бичгийг иш татсан. Claude
    ассистантыг гол төлөв debugging, refactoring, DirectML-д нийцтэй
    алгоритмын хэрэгжүүлэлт (FocalLoss-ыг \texttt{torch.gather} ашиглан
    дахин зохион байгуулах гэх мэт) болон grid search pipeline-ыг
    бүтээхэд ашигласан. Gemini-г санааг харьцуулах, сонголт хийхэд
    зөвлөгөө авах зорилгоор нэмэлт хэрэглсэн. Бүх код нь зохиогчийн
    гараар шалгагдаж, туршигдсан бөгөөд эцсийн хувилбарт өөрийн
    хяналтаар багтсан.
  \item \textbf{Бичгийн туслалцаанд (Claude, Gemini):} Дипломын текстийн
    редакцийн үе шатанд (хэл найруулга, бүтцийн зохион байгуулалт,
    LaTeX-ын техникийн тусламж, томьёоны бүтэц) Claude-ыг үндсэн
    зөвлөх хэрэгсэл болгон ашигласан; Gemini-г зарим тодорхой асуултад
    хариулт өгөх, эх сурвалж хайх зорилгоор ашигласан. Судалгааны
    санаа, арга зүйн сонголт, үр дүнгийн шинжилгээ, дүгнэлт болон
    чухал шийдвэрүүд нь бүхэлдээ зохиогчийн оруулсан бодит хувь нэмэр
    юм. AI хэрэгслийн гаргасан санал дүгнэлтийг шууд хуулах биш, шүүн
    үзэж, зохиогчийн өөрийн үгээр дахин боловсруулан оруулав.
  \item \textbf{Өгөгдлийн цуглуулалт ба шошголт:} Ямар ч AI хэрэгсэлд
    шошгын шийдвэр хийлгээгүй; бүх 4 ангиллын шошго нь хүний
    annotator-ын гараар хийгдсэн.
  \item \textbf{Эцсийн загвар:} MN-BERT-ийн fine-tuning, baseline-уудын
    сургалт, Focal Loss + cosine scheduler бүхий grid search, мөн
    туршилтын үнэлгээ нь бүгд зохиогчийн локал орчинд (AMD RX 9070 GRE
    12GB, \texttt{torch-directml}) явагдсан. AI ассистант нь загварын
    жин, гиперпараметрын эцсийн сонголт эсвэл тайлан бүртгэлд шууд
    нөлөөлөөгүй --- зөвхөн debugging, код оновчлолд зөвлөгч үүрэгтэй
    оролцсон.
\end{itemize}

Энэхүү ил тод мэдүүлэлт нь академик үнэнч байдлыг (academic integrity)
хадгалахад чухал бөгөөд guideline-ийн зүйл 7-ын шаардлагыг бүрэн хангаж байна.

%-------------------------------------------------------------------------------
\section{Бүлгийн дүгнэлт}

Энэ бүлэгт хоёр шатлалт ангиллын системийн ерөнхий архитектур, Stage~1 болон
Stage~2 загваруудын зохион байгуулалт, MN-BERT fine-tuning-ийн стратеги,
суурь загваруудын тохиргоо, тэнцвэргүй өгөгдлийн шийдлүүд, Gradio
интерфейсийн бүтэц, болон үнэлгээний арга зүйг тодорхойлсон. Дараагийн
бүлэгт эдгээр загваруудыг бодитоор сургаж, туршилтын үр дүнг танилцуулна.
